
This work frames blink detection in epoch-structured EEG as both a detection problem and a
threshold-selection problem. Using a shared event-level evaluation protocol, we compared
five strategies and performed targeted stability analyses to understand how methodological
choices (epoch length and event-matching strictness) affect reported performance.

Under the primary Proposed-Med configuration, a 30-second epoch duration achieved the best
macro-F1 among the tested durations (20, 30, 40, 60, 120\,s), with performance remaining
stable across this range (macro-F1\,=\,0.8177 at 30\,s). Strategy comparisons showed
that the proposed epoch-aware pipeline substantially improves over naive concatenation and
a community baseline (MNE annotate\_amplitude), while a direct Bayesian-optimisation
ablation can remain competitive in macro-F1. Finally, the IoU sweep highlights that
evaluation strictness must be reported explicitly: stricter overlap thresholds lead to
markedly lower F1 even when detections are near-misses.